\chapter{One-week optical systems interview review}
\label{app:interview-review}

Use this appendix as a standalone review sheet for about a week of focused
prep. Its purpose is not to cover the most optics. Its purpose is to make your
engineering process automatic under time pressure.

\keyidea{Every answer should end with the engineering decision. Interviewers
remember decisions more than measurements. Close with ``Therefore I would
continue validation,'' ``stop shipment,'' ``contain Supplier~B,'' or ``update
the ATP.''}

\section{Three principles}
\label{sec:interview-philosophy}

\noindent\textbf{Principle 1: Engineering reduces uncertainty.}
\keyidea{The purpose of engineering is not to find certainty. It is to reduce
uncertainty enough to make the next decision.}
\keyidea{The job is not finding the truth. The job is making the best decision
with today's evidence.}
Validation, measurement, debugging, qualification, supplier choices, and
production are the same work under different names. Ask the scale of the
problem (device, module, rack, fleet) before you chase a root cause: scale
picks the owner.

\noindent\textbf{Principle 2: Measurements exist to unlock decisions.}\\
Instruments do not exist to produce plots. They exist to unlock an action.
A power meter asks whether you should chase the optical path or signal
integrity. An OSA asks whether spectral alignment is still plausible. An LIV
asks whether the device itself changed. A DCA\termnote{DCA} or
TDECQ\termnote{TDECQ} asks whether the eye is still inside budget. A BER
waterfall asks whether you have a sensitivity shift or a noise floor. On every
answer, name the decision unlocked (ship, derate, second-source,
ATP\termnote{ATP} change, partner action) as well as the instrument
(\Cref{tab:interview-decisions}).

\noindent\textbf{Principle 3: Measurements characterize margin.}\\
Engineering is not only proving that a product works. It determines how much
uncertainty and margin remain before failure. Debugging identifies exhausted
margin. Qualification verifies that remaining margin is still acceptable after
expected stresses. The same ledgers (power, noise, timing, spectral, control)
appear in both jobs.

\noindent\textbf{Principle 4: Every measurement updates your beliefs.}\\
Treat engineering as hypothesis testing:
\[
\begin{split}
\text{observation} &\longrightarrow \text{hypotheses}
\longrightarrow \text{measurement}\\
&\longrightarrow \text{belief updated / hypothesis eliminated}.
\end{split}
\]
\keyidea{Debugging is simply Bayesian inference performed in a laboratory.}
Every measurement updates the probability of competing hypotheses. Speak the
update out loud. ``At this point I am about 70\% on calibration drift, 20\% on
wavelength walk, and 10\% on receiver noise; before I change firmware I would
verify the eye with a bias sweep.'' A test that leaves the weights unchanged
was the wrong test, or the wrong reference plane.

\paragraph{Engineering priors.}
Before touching the bench, assign higher probability to common failure modes
than to exotic ones. Calibration drift is more likely than simultaneous laser
aging and receiver degradation. A supplier-specific hot-corner escape is more
likely than a new physics mechanism. Choose the first measurements to test
those higher-probability hypotheses while eliminating as many alternatives as
possible. Priors are not prejudice; they are how you spend lab hours.

\keyidea{Engineering is decision making. Decision making is uncertainty
reduction. Measurements reduce uncertainty. Therefore measurements exist to
improve decisions.}

This role owns laser direction inside an IM/DD interconnect effort. Lab
measurement is how you decide, not the whole job. Hands-on fluency still
matters: LIV, RIN, ORL, TDECQ, and a BER waterfall. Senior people lose the
level if they stop at plots and never close on a decision and a control.

Work the day plan at the end of this appendix. Memorize the one-page cheat
sheet before Day~7. Night-before drill is \Cref{app:interview-frameworks}:
open the matching thirty-second framework. Use each LLM practice box the day
it is scheduled. Do not add new topics after Day~6 beyond that drill.

\section{Engineering decision trees}
\label{sec:interview-decision-trees}

Interview questions are usually solved by walking a sequence of
uncertainty-reduction decisions. The exact branches differ. The philosophy
never changes. Debugging asks what broke and which margin ledger is spent.
Qualification asks what uncertainty remains before shipment. The playbooks in
\Cref{app:interview-frameworks} are specialized trees under these two
universal ones.

\paragraph{Universal Tree 1: Debugging.}
\begin{dectree}
Problem
  |
Scope (unit -> lot -> vendor -> fleet)
  |
Population / time behavior
  |
Power changed?
  |-- YES --> Power ledger --> optical path
  |-- NO  --> Quality / receiver
  |
Isolation
  |
Root cause
  |
Decision
  |
Recurrence control
\end{dectree}

\paragraph{Universal Tree 2: Qualification.}
\begin{dectree}
New product
  |
Works? (bring-up)
  |
Margins?
  |
Environmental?
  |
Interoperability?
  |
Reliability?
  |
Manufacturing / ATP?
  |
Pilot deployment?
  |
Production
  |
Fleet monitoring
\end{dectree}

\section{The answer spine}

Move every answer through the same sequence. Memorize four phases, not nine
nodes:
\begin{dectree}
Understand:   requirements -> architecture
Investigate:  measure -> observe -> hypothesize -> isolate
Resolve:      root cause -> corrective action
Prevent:      recurrence control -> decision
\end{dectree}
The diagram is a memory aid, not the answer. Speak one clear paragraph per
phase under time pressure, and expand a node only when asked. Do not jump from
a symptom to a component. End every debug answer with the decision
(\Cref{tab:interview-decisions}). The systems loop in
\Cref{sec:systems-engineering-loop}, the debugging pyramid in
\Cref{sec:debug-pyramid}, and the failure-analysis method in
\Cref{ch:failure-modes} are the full versions of this spine.

\paragraph{Requirements.}
Start by stating what the system must do before you name a part. Reach, lane
rate, aggregate capacity, power envelope, lifetime, cost, and manufacturing
volume are the constraints that decide everything downstream. A 500~m DR
link\termnote{DR} at 100~G/lane with a tight power budget is a different
problem from a 2~km FR\termnote{FR} link or a dense WDM\termnote{WDM}
co-packaged engine. Say the requirement out loud even when the interviewer
hands you a failed module, because the requirement tells you which margins
matter and which architectures were even eligible. If you cannot name the
requirement, you cannot tell whether a measurement is relevant.

\paragraph{Architecture.}
Translate the requirement into an architecture path: fiber plant, wavelength,
source, modulator, receiver, and digital stack. A VCSEL\termnote{VCSEL} path
commits you to 850~nm multimode fiber and direct modulation. A
DFB\termnote{DFB} or EML\termnote{EML} path at 1310~nm commits you to
single-mode fiber and leaves a choice among direct modulation, electro-
absorption, a silicon MZM\termnote{MZM}, or a ring\termnote{MRM}. Each path
then sets detector material, thermal control, test coverage, and service
policy. Architecture is where you show that component choices are not
preferences; they are consequences of the requirement
(\Cref{tab:source-mod-matrix}).

\paragraph{Measurements.}
Name the instruments you would use and the uncertainty each one removes.
Always state the reference plane, pattern, temperature, and pass criterion
with the instrument. A number without a plane is not a measurement. Ask out
loud: what decision does this measurement unlock?

A power meter does not ``measure power'' as an end in itself. It answers
whether the optical path still owns the failure, or whether you should move to
signal integrity. An LIV\termnote{LIV} setup answers whether the laser itself
changed threshold or slope. An OSA\termnote{OSA} answers whether spectral
alignment is still plausible (wavelength and SMSR\termnote{SMSR}).

A DCA\termnote{DCA} answers whether the eye is still inside budget:
OMA\termnote{OMA}, ER\termnote{ER}, RLM\termnote{RLM},
TDECQ\termnote{TDECQ}, and eye shape. A BERT\termnote{BERT} and
FEC\termnote{FEC} counters answer whether you have a sensitivity shift, a
floor, or a burst pattern.

A VNA\termnote{VNA} answers electrical or electro-optic bandwidth. An
ORL\termnote{ORL} meter and inspection scope answer whether the plant is
reflecting or dirty.

\paragraph{Observations.}
Report what the instruments actually showed, not what you hoped they would
show. Separate facts from interpretation. ``Received power held at
$-4$~dBm while pre-FEC BER rose from $10^{-8}$ to $10^{-5}$ at $70^\circ$C''
is an observation. ``The laser is dying'' is not. First scope the failure
(unit $\rightarrow$ lot $\rightarrow$ vendor $\rightarrow$ fleet), and note
sudden versus gradual and recoverable versus permanent. Preserve telemetry
before you reseat, clean, reboot, or rewrite a calibration table.
Observations that disappear under debugging are observations you never had.

\paragraph{Hypotheses.}
Turn observations into a short ranked list. Start from engineering priors
(common modes first), then apply the power-versus-quality fork
(\Cref{sec:interview-power-fork}). Keep the list short enough that the next
measurement can kill more than one item. A hypothesis you cannot falsify with
a bench step does not belong on the list yet.

\paragraph{Isolation.}
\keyidea{Great engineers perform the minimum measurement that eliminates the
largest number of hypotheses.}
Good engineers perform measurements. Do not perform two experiments when one
can separate the hypotheses. Optimize uncertainty removed per hour of lab
time, not uncertainty removed in the abstract. A fast sequence often beats a
complete one:
\[
\begin{split}
\text{golden swap} &\longrightarrow \text{power meter}
\longrightarrow \text{bias sweep}\\
&\longrightarrow \text{then DCA / OSA / RIN as needed}.
\end{split}
\]
A golden swap of transmitter versus receiver splits Tx from Rx. An
optical-versus-electrical lane swap in a multi-lane module splits the optical
path from the driver or TIA\termnote{TIA}. A bias sweep at the failing
temperature asks whether the optimum moved away from the stored table. An
LIV\termnote{LIV} compared with ship data asks whether the device aged or only
the setpoint drifted. A BER-versus-power waterfall asks whether the curve
shifted or floored.

\paragraph{Root cause.}
State the mechanism that survived isolation, with the evidence that killed
the alternatives. ``EAM\termnote{EAM} bias table segment wrong above
$60^\circ$C'' is a root cause. ``Bad eye at high temperature'' is still a
symptom. Name the physical or process mechanism when you can: facet wear,
monitor-PD corruption, FAU\termnote{FAU} misalignment, MPI\termnote{MPI}
from a dirty connector pair, control ledger exhausted (TEC\termnote{TEC} or
heater at rail), supplier lot with high threshold. If the evidence only
reaches ``calibration drift'' and not a deeper mechanism, say so.
Overclaiming the root cause is worse than stopping one layer early with
honesty. Sometimes the product decision is due before the mechanism is known;
that case is \Cref{sec:interview-under-uncertainty}.

\paragraph{Corrective action.}
Fix the mechanism you named, under the condition that failed. Retune the
table, replace the lot, change the thermal design, clean and inspect the
plant, filter the supply, or rewrite the ATP\termnote{ATP} limit. Containment
comes first when the fleet is already exposed: stop ship, quarantine lots, or
derate while the permanent fix is built. Then walk the verification ladder:
reproduce the failure, verify the fix, regression-test neighbors, requalify if
the change touches life or safety, release to production, and watch fleet
telemetry. Candidates often stop at ``I fixed it.'' Staff engineers do not.

\paragraph{Recurrence control and the decision.}
Close with the control that stops the same escape next month, then state the
product decision in one sentence. The control is usually a new or tightened
test, an SPC\termnote{SPC} chart, a telemetry alarm, a process-control limit,
a supplier screen, or a firmware guard that refuses to boot with a railed
actuator. Name the owner and the measurable closure criterion. Then say the
action: continue validation, stop shipment, contain the lot, update ATP,
escalate the supplier, or monitor only. Root cause without recurrence control
is a story. Recurrence control without a decision is incomplete. In the
interview, ending on the decision is what separates a debug narrative from a
product-engineering answer.

\keyidea{I first want to understand the scope of the problem, then determine
which margin ledger is being spent, choose the measurement that eliminates the
largest number of hypotheses, make the product decision, and finally add the
control that prevents the next escape.}

\llmpractice{Four principles: uncertainty reduction; measurements unlock
decisions; measurements characterize margin; every measurement updates belief.
Chunk the spine as Understand / Investigate / Resolve / Prevent. Every answer
ends with a decision. Name the ledger that moved.}
{Quiz me on the four principles, priors, and the four spine phases. Give me a
failed module at high temperature with stable average power. Stop me if I skip
scope, the power fork, the control ledger, the product decision, or recurrence
control.}

\section{Staff judgment under uncertainty}
\label{sec:interview-under-uncertainty}

The spine above assumes you can reach a mechanism. Senior work often cannot.
The product decision still has a deadline. Staff engineers make the decision
with the evidence they have, name residual risk, and keep the FA path open.

\subsection{Decide before the mechanism is known}

Evidence can be enough even when the physical story is not. Example:
\begin{description}
\item[Evidence] BER degrades. Only Supplier~B. Only the hot corner. About
      3\% of the population.
\item[Mechanism] Unknown.
\item[Decision] Stop shipment of Supplier~B. Continue Supplier~A. Expand the
      ATP\termnote{ATP} hot-corner screen. Begin FA\termnote{FA} and request
      DPA\termnote{DPA} on fails. Review telemetry and RMA\termnote{RMA} daily
      until the mechanism closes.
\end{description}
That is real engineering. Waiting for SEM photos while bad lots keep shipping
is not.

\subsection{Measurements unlock actions}

Measurements do not unlock understanding as an end in itself. They unlock
actions. Keep this vocabulary ready and use it in answers
(\Cref{tab:interview-decisions}).

\begin{table*}[ht]
\refstepcounter{table}\label{tab:interview-decisions}%
\centering
\setlength{\tabcolsep}{4pt}%
\footnotesize
\begin{tabular}{@{}p{0.22\linewidth}p{0.72\linewidth}@{}}
\toprule
Action & Typical unlock \\
\midrule
Ship / don't ship & Population meets ATP and life model, or it does not \\
Continue validation & Uncertainty still blocks a ship or architecture call \\
Escalate supplier & Lot, site, or vendor signature after scope \\
Derate & Margin too thin; product can ship under tighter use \\
Second source & Single-vendor risk exceeds fleet tolerance \\
Contain lot & Date-code or lot escape; stop further exposure \\
Modify ATP & Escape path found; production must catch it next \\
Open RMA / request FA & Field or partner unit needs mechanism work \\
Perform DPA & Need physical confirmation of facet, solder, FAU, die \\
Change firmware & Control loop, table, or guard is wrong \\
Retune calibration & Device healthy; setpoint or table segment wrong \\
Monitor only & Rate tiny, flat, no customer impact; watch trends \\
\bottomrule
\end{tabular}
\par\medskip
{\raggedright\footnotesize\textbf{Table~\thetable.} Decision vocabulary for
interview answers. Name the action, the owner, and the residual risk when the
mechanism is still open.\par}
\end{table*}

\subsection{Time is a resource}

Every measurement costs hours. Optimize uncertainty removed per hour of lab
time. Prefer the cheap cut that kills many hypotheses before the slow
characterization. Telemetry and a golden swap often beat an immediate RIN
setup. SEM and DPA come last, not first.

\subsection{Measurement hierarchy}

Not every failure deserves destructive analysis. Climb only as far as the
decision requires:
\begin{dectree}
telemetry -> simple bench -> swap -> characterization -> FA / DPA
\end{dectree}

\subsection{Fleet economics}

Product engineering is economics as well as physics. A tiny, flat, no-impact
fail rate can stay on a monitor-only plan. A growing, supplier-specific rate
demands containment the same day.
\begin{description}
\item[Monitor only] Failure rate $\sim$0.003\%, no customer impact, no trend,
      no growth. Keep telemetry alarms and review weekly.
\item[Contain immediately] Failure rate $\sim$2\%, growing, supplier-specific.
      Stop ship, quarantine lots, notify the partner, open FA, tighten ATP.
\end{description}
Intentionally not fixing a root cause is sometimes correct. Leaving a growing
escape unowned is never correct.

\subsection{Ownership language}

When the interviewer asks ``What would you do?'', answer as the owner:
\begin{quote}
As owner I would stop shipment of the affected lots, notify the supplier,
request FA with DPA on a sample, update the ATP hot-corner screen, review
fleet telemetry and RMA codes daily, and schedule a qualification rerun
before open volume resumes.
\end{quote}
Ownership is the difference between a debug narrative and a Staff answer.

\llmpractice{Staff judgment: decide with incomplete mechanism evidence; name
an action from the decision table; spend lab hours for uncertainty removed;
climb the measurement hierarchy only as far as needed; apply fleet economics;
speak as the owner.}
{Give me three incomplete-evidence scenarios. For each I must state evidence,
confidence weights, the action I take today, residual risk, and the FA path.
Fail me if I wait for a perfect root cause before containing a growing lot.}

\section{Common interview traps}
\label{sec:interview-traps}

Interviewers listen for these mistakes. Avoid them on purpose.

\begin{description}
\item[Trap 1: Naming a component first.] Start with scope, not a part number.
\item[Trap 2: Calling symptoms root causes.] ``Bad eye'' is a symptom.
      ``EAM bias table wrong above $60^\circ$C'' is a root cause.
\item[Trap 3: Stopping after the fix.] Always end with recurrence control and
      the product decision.
\item[Trap 4: Listing instruments.] Name the decision each measurement unlocks,
      not a gear catalog.
\item[Trap 5: One unit as the fleet.] Scope to lot, vendor, rack, and trend
      before you generalize.
\end{description}

\keyidea{Before any debugging answer, ask: Scope? Which ledger moved? Power or
quality? Fastest measurement? Decision? Control?}

\section{Ten concepts to know cold}

\subsection{Start at the system}

For a new link, walk downward from capacity to the service model. Capacity
sets lane rate. Lane rate and reach together set the fiber plant and the
source class. Power and cost then decide whether you can afford a
TEC\termnote{TEC}, a retimer, or a dense WDM\termnote{WDM} engine. Only after
those constraints are named do you choose the modulator and the receiver,
then the DSP\termnote{dspserdes} and FEC\termnote{FEC} stack, then the
validation and field-service model.
\[
\begin{split}
\text{capacity} &\longrightarrow \text{lane rate}
\longrightarrow \text{reach}
\longrightarrow \text{power}\\
&\longrightarrow \text{fiber plant}
\longrightarrow \text{source}
\longrightarrow \text{modulator}\\
&\longrightarrow \text{receiver}
\longrightarrow \text{digital signal processing and FEC}\\
&\longrightarrow \text{validation and service model}.
\end{split}
\]

A good answer explains why each choice constrains the next one. A
VCSEL\termnote{VCSEL} path points toward 850~nm, multimode
fiber\termnote{MMF}, silicon detection, and direct modulation with a short
reach. A DFB\termnote{DFB} path at 1310~nm points toward single-mode
fiber\termnote{SMF}, germanium detection, and a choice among a directly
modulated laser\termnote{DML}, an EML\termnote{EML}, or an external
silicon-photonic modulator. Neither path is ``better.'' Each closes a
different reach, power, cost, and manufacturing problem. The decision matrix
is in \Cref{tab:source-mod-matrix}.

\subsection{Scope before root cause}

Before you open an instrument, walk the failure up the scope ladder. Each rung
changes the owner and the next action:
\[
\begin{split}
\text{failure observed} &\longrightarrow \text{single unit?}
\longrightarrow \text{lot?}\\
&\longrightarrow \text{vendor?}
\longrightarrow \text{rack?}\\
&\longrightarrow \text{datacenter?}
\longrightarrow \text{fleet?}
\end{split}
\]
Also ask time and change: new failure or long-running trend; sudden, gradual,
intermittent, or temperature-dependent; firmware, supplier, process, fixture,
workload, or environment change just before the symptom. Scope often removes
more hypotheses than the first bench measurement. A fleet-wide gradual drift
cannot be a single dirty connector. A single-lane sudden failure after a rework
cannot be a population aging problem. A vendor-lot signature points to supplier
containment before you redesign the module.

Preserve the failing state and its telemetry before you reseat, clean,
reboot, or change calibration. Capture CMIS\termnote{CMIS} monitors, pre-
FEC\termnote{FEC} BER\termnote{BER}, bias currents, temperatures,
LOS\termnote{LOS} and LOL\termnote{LOL} flags, and firmware versions. An
intermittent that disappears under debugging is still a real failure; you
just destroyed the evidence
(\Cref{sec:fleet-triage,tab:failure-analysis-checklist}).

\subsection{Use the power-versus-signal-quality fork}
\label{sec:interview-power-fork}

First ask whether received optical power changed. That one question splits
the debug tree.

If power changed, the power ledger moved. Inspect source output and enable
state, coupling and FAU\termnote{FAU} alignment, connector cleanliness,
ORL\termnote{ORL}, fiber and multiplexer loss, and monitor-photodiode
calibration. An APC\termnote{APC} loop that trusts a drifting monitor will
hold a wrong launch power while reporting that everything is fine. Confirm
with an external power meter at a named reference plane.

If power held but BER worsened, the loss is in signal quality or in the
receiver. Inspect OMA, ER, RLM, and TDECQ\termnote{TDECQ} on a DCA;
RIN\termnote{RIN} under controlled ORL; jitter and ISI; wavelength alignment
to filters or rings; receiver sensitivity and equalization; and calibration
tables that set modulator bias. Also ask the receiver-side equivalent: did
the required optical power increase even though the power meter reading did
not? A sensitivity shift can look exactly like transmitter degradation until
you do a golden swap (\Cref{sec:debug-fork,sec:validation-fork}).

\subsection{Track five margin ledgers}
\label{sec:interview-margin-ledgers}

Links rarely fail from one dramatic excursion. They fail when several small
shifts spend different ledgers at once. Almost every failure can be described
as spending one or more ledgers. Use that language in every debug answer.
Name which ledger moved, by how much, and what spent it.
\begin{dectree}
Power · Noise · Timing · Spectral · Control
\end{dectree}

\begin{description}
\item[Power] Launch power, insertion loss, coupling, connector loss,
      multiplexer loss, and receiver sensitivity. Track average power and
      OMA\termnote{OMA} separately: average power can look healthy while OMA
      collapses.
\item[Noise] RIN\termnote{RIN}, receiver thermal noise, shot noise,
      crosstalk, supply noise through weak PSRR, and MPI from reflective
      interfaces. Noise that scales with the signal makes a BER floor that
      more launch power cannot fix.
\item[Timing] Bandwidth, dispersion, jitter, ISI, and equalization reserve.
      A SerDes\termnote{SerDes} that is already using most of its FFE and DFE
      taps has little timing margin left even if the eye still opens on a DCA.
\item[Spectral] Wavelength, SMSR\termnote{SMSR}, filter or ring passband,
      thermal drift, and lock range. A ring whose heater is near its DAC rail
      is spending spectral margin even while BER still passes.
\item[Control] Headroom in the loops that keep the optics healthy:
      APC\termnote{APC}, TEC\termnote{TEC}, heaters, ring lock, bias DACs, and
      calibration tables. Prefer product language over bench language: say
      ``the control ledger is exhausted'' or ``we have no remaining control
      authority,'' not only ``TEC current hit maximum'' or ``heater DAC went
      full scale.'' A railed actuator or a corrupted table can fail the link
      while the laser diode itself is still healthy. Control margin is what
      separates ``optics broke'' from ``the product can no longer hold the
      operating point.''
\end{description}

The device and link treatment is in \Cref{sec:laser-margin-erosion}. Always
ask the control ledger when actuators and tables are in the product.

\subsection{Use the validation ladder}
\label{sec:interview-validation-ladder}

\keyidea{Validation is staged uncertainty reduction. Engineering is decision
making; decision making is uncertainty reduction; measurements reduce
uncertainty; therefore measurements exist to improve decisions.}

For every stage, name the question the stage answers and the uncertainty it
removes. A test that answers no question is cost, not confidence. Later
sections only say ``walk the validation ladder''; this is the full map.

\begin{description}
\item[Bring-up] Does the hardware power, initialize, emit, lock, and pass
      data on a known-good host? Removes integration unknowns only.
\item[Characterization] How does performance vary with temperature,
      voltage, lane, wavelength, pattern, and unit? Maps the population, not
      a hero sample. Record LIV, spectrum, RIN, OMA, ER, RLM, and
      TDECQ\termnote{TDECQ} as distributions.
\item[Margin testing] How far is the design from each failure boundary?
      Push power, temperature, ORL\termnote{ORL}, and supplies to the cliff
      and measure the distance in dB, degrees, and volts.
\item[Interoperability] Does it work with target hosts, production fiber,
      real modules, and real firmware? Removes combination failures that
      golden benches miss.
\item[Environmental testing] What changes under temperature cycling,
      humidity, vibration, shock, and connector mating?
\item[Reliability qualification] Which parameters drift with stress and
      time? HTOL\termnote{HTOL}, with a named mechanism and justified
      activation energy, projects wear-out as a FIT\termnote{FIT} number.
\item[Production readiness] Can the ATP\termnote{ATP} catch bad units and
      process drift at useful cost and cycle time, with station correlation
      and guard bands sized to repeatability?
\item[Fleet monitoring] Which telemetry catches margin erosion before
      service failure, and does the RMA\termnote{RMA} Pareto match the life
      model?
\end{description}

\Cref{tab:ladder} maps these stages to activities and instruments. The
worked answer in \Cref{sec:interview-worked-validation} is practice prose
built on this ladder. Night-before path: Framework~7 in
\Cref{app:interview-frameworks}.

\subsection{Margin budgeting}
\label{sec:interview-margin-budget}

Every environmental or use stress consumes part of the system margin:
temperature, voltage variation, supply ripple, fiber contamination, insertion
loss, connector wear, aging, mechanical vibration, and process variation.
Qualification is not a tour of each mechanism for its own sake. It verifies
that after the expected stresses, remaining margin is still acceptable.
Stress consumes margin. Qualification measures remaining margin. Debugging
finds which ledger is exhausted when that remaining margin hits zero
(\Cref{sec:interview-margin-ledgers}).

\subsection{Customer view versus vendor view}
\label{sec:interview-customer-vendor}

The vendor designs internals. The customer characterizes externally
observable behavior. As a customer you often do not need laser threshold,
driver architecture, or TIA topology. You measure BER, sensitivity, FEC
statistics, launch and receive power, telemetry, and environmental response;
eye metrics when engineering access exists. If the product is a black box,
qualification focuses on that external surface. If engineering samples are
available, request transmitter-only, receiver-only, breakout, or diagnostic
hardware to isolate Tx and Rx margins independently. Keep the view explicit
in second-source and qualification answers
(\Cref{sec:interview-worked-second-source}).

\subsection{Know what each instrument answers}

Do not recite instrument names. Say what uncertainty each one removes, and
what decision that unlocks.

\begin{description}
\item[Power meter] Did optical power change at a named reference plane?
      Decision: chase the optical path, or move to signal integrity.
      External meters catch monitor-PD lies.
\item[LIV setup] Did laser\termnote{LIV} threshold, slope efficiency,
      voltage, kink behavior, or thermal rollover change relative to ship
      data? Decision: has the device itself changed?
\item[Optical spectrum analyzer or wavemeter] Did wavelength,
      SMSR\termnote{SMSR}, mode structure, or spectral width change on an
      OSA? Decision: is spectral alignment still plausible for filters and
      rings?
\item[RIN setup] Is transmitter intensity noise\termnote{RIN} limiting BER,
      and does it depend on ORL or on the product bias board rather than the
      laser itself? Decision: is the floor optical noise or something else?
\item[Digital communication analyzer] Are OMA, ER, RLM,
      TDECQ\termnote{TDECQ}, jitter, and eye shape acceptable on a DCA?
      Decision: is the eye still inside the transmitter budget at this corner?
\item[Bit-error-ratio tester and FEC counters] What is the BER waterfall,
      floor, burst pattern, lane correlation, and dwell-time behavior from a
      BERT\termnote{BERT}? Decision: sensitivity shift, noise floor, or
      intermittent? FEC histograms separate Gaussian noise from clustered MPI.
\item[Vector network analyzer] Is electrical or electro-optic bandwidth,
      impedance, or reflection causing eye closure on a VNA\termnote{VNA}?
      Decision: is the plant electrical rather than optical?
\item[ORL meter and inspection scope] Is the optical plant adding loss or
      reflection\termnote{ORL}? Decision: clean and inspect before blaming the
      laser.
\item[Thermal chamber] Is behavior reversible with temperature, and which
      actuator (TEC\termnote{TEC}, heater, bias DAC) loses headroom?
      Decision: thermal operating point versus aging.
\item[Bias sweep] Where is the optimum for the EAM\termnote{EAM}, MZM, or
      ring, and did it move away from the stored table? Decision: retune the
      control ledger, or replace the device?
\end{description}

Always name the reference plane, setup calibration, test pattern, bandwidth,
temperature, sample identity, and pass criterion with the result
(\Cref{tab:measurement-mapping,tab:laser-meas}).

\subsection{Read a BER waterfall: shift, floor, and burst pattern}
\label{sec:interview-waterfall}

These three words show up in almost every debug answer. Know what each one
looks like on the bench, what it rules in, and what it rules out. The
operational procedures are in \Cref{sec:fm-ber-increase,sec:fm-ber-floor}.

\paragraph{What a BER waterfall is.}
A BER waterfall is a plot of bit error ratio versus received optical power.
You build it by sweeping a calibrated VOA\termnote{VOA} in the path, holding
pattern, temperature, and host fixed, and counting errors long enough at each
power that the BER\termnote{BER} is statistically meaningful. Plot received
power on the horizontal axis (name the reference plane: usually TP3 at the
receiver connector) and pre-FEC BER on a log vertical axis. A healthy link
falls steeply as power rises: more photons, better signal-to-noise ratio,
fewer errors. That falling curve is the waterfall. A single BER point at the
operating power is not a waterfall. Without the sweep you cannot tell whether
you are short on power or limited by noise that scales with the signal.

\paragraph{What a shifted waterfall means.}
A parallel shift means the whole curve moved left or right while keeping a
similar slope. The link still improves when you add power; it just needs a
different power to hit the same BER. A rightward shift (worse sensitivity)
means you now need more received power for the same pre-FEC BER. Common causes
are lost launch or coupling (power ledger), a quieter or noisier receiver
(TIA\termnote{TIA} noise, responsivity), eye closure that lowers effective
OMA\termnote{OMA} at constant average power (ER\termnote{ER}, RLM\termnote{RLM},
TDECQ\termnote{TDECQ}), wavelength walking onto a filter edge, or equalizer
misadaptation. A leftward shift means the link got healthier. The diagnostic
move after you see a shift is a golden swap: known-good transmitter, then
known-good receiver, to decide which side owns the sensitivity change. Raising
launch power can still help a shifted link, because the waterfall has not
stopped responding to power.

\paragraph{What a BER floor means.}
A floor is a horizontal asymptote: as you raise received power, BER improves
for a while and then stops. More power no longer helps. That happens when the
dominant noise scales with the signal itself. Relative intensity
noise\termnote{RIN} is the textbook case: intensity fluctuations grow in
proportion to optical power, so the signal-to-noise ratio saturates and
\[
Q_{\max}=1/\sqrt{\mathrm{RIN}\cdot\mathrm{BW}}.
\]
Multipath interference\termnote{MPI}, bias-rail noise through weak
PSRR\termnote{PSRR}, and aggressor crosstalk that tracks the aggressor power
make the same shape. The interview trap is to keep raising launch power or
cleaning for loss when the curve has already floored. The fix is to remove the
multiplicative noise source (ORL\termnote{ORL} and isolator hygiene, supply
filtering, connector reflections, crosstalk isolation), not to buy more
photons. Confirm the floor with a full sweep before you name the mechanism;
then bisect with a quiet laser source versus the product bias board, an ORL
reflector sweep, and the FEC error timing below.

\paragraph{What a burst pattern means.}
Average BER alone hides how errors arrive in time. A burst pattern means
errors cluster: many errored symbols in a short window, then quiet intervals,
rather than a steady sprinkle of random bit flips. FEC\termnote{FEC}
histograms and pre-FEC error counters with timestamps make this visible.
Gaussian thermal noise and well-behaved RIN tend to spread errors. MPI from a
pair of reflective interfaces, connector intermittents, ESD events, supply
glitches, and unlocked CDR\termnote{CDR} intervals tend to cluster them. Lane-
correlated bursts across a module point at a shared supply, clock, or thermal
event. A single-lane burst pattern points at that lane's optical path or
connector. In the fleet, bursty pre-FEC counters with stable average power are
often intermittents: preserve the counters and CMIS\termnote{CMIS} event log
before you reseat anything, or the evidence disappears
(\Cref{sec:fm-intermittent}).

\paragraph{How to say the three together.}
``I sweep BER versus received power. A parallel shift is a sensitivity or
power-ledger problem; more power can still help. A floor is multiplicative
noise; more power cannot help. A bursty FEC histogram means clustered
impairments such as MPI or intermittents, not plain Gaussian noise.'' Offer
to draw the shifted curve next to the floored curve on the whiteboard.

\subsection{Separate thermal response from aging}

Thermal effects are usually reversible when you return to the starting
temperature: wavelength shift, ring detuning, EAM\termnote{EAM} or
MZM\termnote{MZM} bias movement, TEC\termnote{TEC} loading, and
receiver-noise increase. Aging is cumulative and does not reverse on
cool-down: threshold rise, slope loss, contact degradation, defect growth,
and permanent absorption or spectral change.

The practical test is simple. Return to the starting temperature and compare
with pre-stress data at the same junction temperature. Full recovery supports
a thermal or control hypothesis, often a calibration-table segment or an
actuator that ran out of range. A permanently moved LIV\termnote{LIV}
baseline supports aging, damage, or assembly change. Mixing the two owners
wastes weeks: a reliability team cannot fix a wrong table, and a firmware
team cannot fix a worn facet (\Cref{sec:laser-thermal-vs-aging}).

\subsection{Accelerated life tests need a mechanism}

HTOL\termnote{HTOL} predicts field life only if elevated temperature and
bias accelerate the same physical mechanism the fleet will see. The stress
must not introduce a different mechanism, such as solder creep at a
temperature the product never reaches, and it must not omit a dominant field
stress, such as thermal cycling or connector wear.

In the interview, state the assumed mechanism, the activation energy and why
it applies to this process, the sample size and confidence bound, the stress
and use temperatures, the measured drift parameter (threshold, slope,
wavelength), and the field-data check. Treat the projected FIT\termnote{FIT}
result as a hypothesis that field returns must confirm or falsify. A FIT
number without those pieces is a spreadsheet, not a prediction. The life
acceleration model is Arrhenius\termnote{Arrhenius}
(\Cref{sec:laser-aging}).

\subsection{Calibration is part of the product}

The operating points a product actually stores are as much the product as the
die. Know them by name: laser bias or APC\termnote{APC} target;
EAM\termnote{EAM} bias versus temperature; MZM\termnote{MZM} quadrature;
ring heater or wavelength-lock point; monitor-photodiode\termnote{PD}
calibration; and receiver and power-monitor offsets. Tables are usually
segmented by temperature, and segment boundaries are a real failure mode.

Recalibrate when evidence demands it, not by habit: loop residual stops
converging, an actuator approaches its rail, telemetry disagrees with an
external reference, temperature leaves the table range, or repair, rework,
or firmware changes invalidate coefficients. ATP\termnote{ATP} must verify
calibration at the temperature corners the fleet will see, not only at
station ambient (\Cref{sec:laser-calibration}).

\subsection{Corrective action must prevent recurrence}

Root cause is not the end of the answer. Close with immediate containment;
the design, process, calibration, firmware, or supplier correction;
verification under the original failing condition; an acceptance test,
process control, alarm, or telemetry control that catches recurrence; and an
owner with a measurable closure criterion. An interview answer that stops at
``we replaced the laser'' sounds like repair. An answer that ends on the new
ATP\termnote{ATP} corner or the new telemetry alarm sounds like product
engineering.

\llmpractice{Ten cold concepts: start at the system; scope before root cause;
power versus signal-quality fork; five margin ledgers; validation ladder;
instruments named by uncertainty removed; waterfall shift versus floor versus
burst; thermal response versus aging; HTOL needs a mechanism; calibration is
part of the product; corrective action must prevent recurrence.}
{Act as my interviewer for optical systems. Ask me one question from each cold
concept, in random order. Push for a named reference plane, a next measurement
that kills hypotheses, and a recurrence control. Grade me on process, not on
jargon volume.}

\section{Two stories to prepare}

Each story is a spoken version of the answer spine. Prepare one component or
bench story and one system, production, or fleet story. Use only work you
personally performed, and separate your contribution from the team's. Walk
the same eight beats every time:

\begin{enumerate}
\item Requirement and system context: what the link or product had to do.
\item Observed symptom and scope: what failed, how wide, sudden or gradual.
\item Competing hypotheses: the short list after the power fork.
\item Measurements chosen and why: which instrument, which uncertainty.
\item Evidence that eliminated each hypothesis.
\item Root cause: the mechanism that survived, with evidence.
\item Corrective action: the fix, verified under the failing condition.
\item Recurrence control and measured result: the test, alarm, or process
      change, and the metric that proved it worked.
\end{enumerate}

\llmpractice{Prepare two true stories, one bench or component and one system,
production, or fleet. Each story must hit the eight beats of the answer spine
and end on a measured recurrence control. Separate your work from the team's.}
{I will tell you my two story titles only. Interview me on each for five
minutes. Force the eight beats in order. Cut me off if I invent metrics, skip
scope, or end without a control that would catch the next escape.}

\section{Questions to rehearse aloud, with model answers}

Rehearse these until the structure is automatic. For night-before review use
the matching playbook in \Cref{app:interview-frameworks}; the prose below is
practice depth. Each answer starts with scope or requirements, names
measurements and the hypotheses they separate, and ends with a decision and a
control. Do not recite; adapt the skeleton to the follow-up questions.

\paragraph{How to use the worked answers.}
Each long question opens with a green \textbf{30-second answer (memorize)}
box. That is what you deliver first. The 3-minute section is for practice.
The 10-minute section is read-only reference. Expand only where the
interviewer asks:
\[
\text{30-second answer}
\longrightarrow \text{interviewer asks}
\longrightarrow \text{expand only there}.
\]

\subsection{How would you set laser requirements for a new IM/DD link?}
\label{sec:interview-worked-laser-reqs}

This is the ownership question. Start from the system, not from a laser
datasheet. The output is a requirements slice a supplier and an ATP can both
test against (\Cref{tab:laser-prd,tab:source-mod-matrix}).

\execanswer{Freeze reach, lane rate, fiber, power, lifetime, and volume.
Those choose the source path. Then write OMA and RIN at a named plane and ORL,
thermal and control headroom, a named HTOL mechanism, and an ATP with supplier
reaction plan. Therefore I would freeze requirements that let us ship and
second-source, not pick the laser that looks best on a bench.}

\paragraph{3-minute answer (practice).}
Walk four steps: (1)~system constraints choose the architecture path before
any part number; (2)~optical budget at named planes (OMA, RIN at ORL, SMSR,
chirp); (3)~thermal, life, and control headroom with a named
HTOL\termnote{HTOL} mechanism; (4)~ATP\termnote{ATP} methods, FAIR triggers,
and RMA\termnote{RMA} codes split by supplier. Name one hard number you would
fight for (RIN under ORL, or APC headroom at hot) and what fails if it is
missing.

\paragraph{10-minute reference (read only).}
Expand one constraint into the full budget table and show how it lands in ATP
limits and partner FAIR triggers. Details: a short multimode path points toward
a VCSEL\termnote{VCSEL}; a 500~m or 2~km single-mode path toward a DFB or
EML\termnote{EML} at 1310~nm. Say the reference plane with every number. If
the laser cannot hold lock or APC headroom at the hot corner, the link does
not close even if room-temperature OMA looks fine.

\subsection{How would you validate a new optical transmitter from bring-up
through production?}
\label{sec:interview-worked-validation}

This is the question most likely to open the interview. The ladder itself is
in \Cref{sec:interview-validation-ladder,tab:ladder}. Frame first: validation
is staged uncertainty reduction. Each stage answers a question the previous
stage could not.

\execanswer{Validation is staged uncertainty reduction. Bring-up proves the
setup is sane. Characterization maps the population. Margin finds the cliff.
Interop removes combination risk. Qual projects life with a named mechanism.
Production proves the ATP catches escapes. Fleet telemetry closes the loop.
Therefore I would walk the ladder in order and refuse any test that answers
no new question.}

\paragraph{3-minute answer (practice).}
Walk the ladder in order. For each stage, name one instrument, the uncertainty
removed, and the decision unlocked (continue, redesign, tighten ATP, stop
ship). End on which ledger the telemetry must watch.

\paragraph{10-minute reference (read only).}
Expand the stage the interviewer picks. The stage-by-stage notes below are
reference detail; do not memorize them.

\paragraph{Stage 1: bring-up.}
For bring-up, get one unit alive on a known-good bench and remove the
integration unknowns. Power the device through its specified sequence and
confirm rails, inrush, and power state transitions. Initialize the firmware and
read identity and status over the management
interface\termnote{CMIS}, because a part that cannot report its own state
cannot be debugged later. Enable the laser and confirm first light on a
calibrated power meter at the specified output reference plane, not at
whatever connector is convenient. Then close a real loop: lock the receiver,
run a PRBS\termnote{PRBS} pattern from a BERT\termnote{BERT} or host SerDes,
and count errors. The uncertainty removed is basic: the design functions, the
chips talk to each other, and the test setup itself is sane. The trap is that
everything here is golden: pristine connectors, one temperature, one host, one
unit. Passing bring-up says nothing about the population, the corners, or the
lifetime. Its only job is to make the next stage's failures meaningful.

\paragraph{Stage 2: characterization.}
For characterization, map what the design actually does across its operating
range, on enough units to see variation. Sweep temperature, supply voltage,
bias, and test pattern, and at each point record the transmitter's vital
signs: the LIV curve\termnote{LIV} for threshold, slope efficiency, and
thermal rollover; the optical spectrum on an OSA\termnote{OSA} for center
wavelength and SMSR\termnote{SMSR}; RIN\termnote{RIN} measured at the stated
ORL\termnote{ORL}, because a quiet-bench RIN number is flattery; and the eye
on a DCA\termnote{DCA} for ER\termnote{ER}, OMA\termnote{OMA}, and
TDECQ\termnote{TDECQ}. The key discipline is to record distributions, not a
golden sample: ten units across two lots tell you what the process gives you,
one hero unit tells you what marketing wants. The uncertainty removed is where
the population sits relative to the spec at every corner. The output is not a
pass stamp; it is the dataset that will later set production test limits and
telemetry alarm thresholds.

\paragraph{Stage 3: margin testing.}
For margin, stop asking whether the part passes and ask how far it sits from
the cliff. Push each parameter to its failure boundary: sweep received power
down with a calibrated VOA\termnote{VOA} until the pre-FEC BER\termnote{preFEC}
waterfall crosses the FEC threshold; take case temperature to the limit and
past it; pull supply rails to their corners; degrade ORL with a controlled
reflector and watch the error floor. Measure the distance to failure in dB,
degrees, and volts, then compare it against the variation found in
characterization. The uncertainty removed is the one that kills fleets: a
design can pass every nominal test while sitting half a decibel from its
limit, and normal lot spread plus aging will spend that half decibel in the
first year. Margin testing also reveals which of the five margins, power,
noise, timing, spectral, or control, runs out first, which tells you what
telemetry must watch (\Cref{sec:interview-margin-ledgers,sec:laser-margin-erosion}).

\paragraph{Stage 4: interoperability.}
For interoperability, take away everything golden. Replace the BERT with the
target hosts, plural, because SerDes equalizers from different vendors adapt
differently to the same transmitter. Replace the lab jumper with production
fiber plant: real connector losses, real MPO\termnote{MPO} polarity, realistic
ORL from the connectors the datacenter will actually use. Run real link
bring-up sequences, firmware interactions, and traffic patterns rather than
one endless PRBS. The uncertainty removed lives in combinations rather than
components: a transmitter and a receiver that each meet spec can still fail
together, because the standard's transmitter eye and stressed-receiver
tests\termnote{SECQ} only guarantee interoperability if both sides were tested
against the specified reference, and reference receivers are idealizations.
Interop failures found here cost a lab week; found in production they cost a
fleet rollback.

\paragraph{Stage 5: environmental stress and reliability qualification.}
For qualification, ask which mechanisms move with stress and time. Each stress
is chosen for the failure mechanism it accelerates, not run as a ritual:
temperature cycling exercises solder joints, die attach, and fiber attach;
damp heat exercises corrosion and delamination; vibration and shock exercise
the mechanical path; HTOL\termnote{HTOL} accelerates the wear-out of the
active region so that weeks of stress project years of field life through an
Arrhenius model\termnote{Arrhenius}. State the projection's fine print
without being asked: the activation energy must be justified for this
mechanism on this process, the sample size sets the confidence interval, and
the projection holds only if the stress accelerates the same mechanism the
field will see. The uncertainty removed is the shape of lifetime: infant
mortality rate for the burn-in decision, and wear-out rate as a
FIT\termnote{FIT} number the fleet math can consume.

\paragraph{Stage 6: production readiness.}
For production, the question changes from is the design good to can the
factory tell a good unit from a bad one, at rate, at acceptable cost. Write
the ATP\termnote{ATP} so that every limit traces to the link budget or to the
characterization data; a limit nobody can trace will be waived under ship
pressure. Correlate stations by running the same golden units on every
station and site, so a limit means the same thing everywhere. Size guard
bands from measurement repeatability, because a test that cannot repeat
tighter than its limit ships escapes and scraps good units in equal measure.
Verify calibration at the temperature corners the fleet will see, not only at
station ambient. Archive raw data, instrument identity, and
calibration-table versions so any shipped result can be replayed. Then keep
SPC\termnote{SPC} charts on threshold current, power, and wavelength, because
process drift between qualification lots is caught here or not at all.

\paragraph{Stage 7: fleet telemetry and the feedback loop.}
Deployment is the last validation stage, not the end of validation. Fleet
telemetry watches the margins that stage 3 identified: per-lane transmit
power, bias current, pre-FEC BER, temperature, and actuator drive such as
TEC\termnote{TEC} current, with alarms on trends and disagreements rather
than only on hard thresholds. Field returns close the loop: the
RMA\termnote{RMA} Pareto is compared against the qualification projection,
and divergence means the life model was wrong, not that the fleet is unlucky.
A high NFF\termnote{NFF} rate is itself a finding, usually pointing at
intermittents or triage gaps. The uncertainty removed is the honest one: what
every earlier stage missed.

\subsection{BER worsens at high temperature but average power is stable. What do
you do?}
\label{sec:interview-worked-hot-ber}

Classic fork question (\Cref{sec:fm-temperature,sec:interview-power-fork}).

\execanswer{Power held, so leave the power ledger. First scope the failure
(unit $\rightarrow$ lot $\rightarrow$ vendor $\rightarrow$ fleet) and whether
cool-down recovers. At the failing temperature, read eye, bias sweep,
wavelength, and control headroom. Therefore I would fix the table or thermal
design and put that loaded corner in the ATP.}

\paragraph{3-minute answer (practice).}
First scope the failure. Apply the power-versus-quality fork: APC is hitting
setpoint, so candidates are ER\termnote{ER} collapse, wrong modulator bias,
wavelength walk, exhausted control ledger (TEC\termnote{TEC} or heater at
rail), or hotter receiver noise. Measure at the failing temperature on a
DCA\termnote{DCA} (ER, OMA, TDECQ), bias-sweep the EAM\termnote{EAM} or MZM,
check OSA wavelength, and read actuator codes. If a bias sweep restores the
eye, retune the table; if the actuator is railed, fix thermal design. Close
with the new ATP corner.

\paragraph{10-minute reference (read only).}
Offer the calibration-table segment-boundary story or the railed-heater story
if asked. Aging does not reverse on cool-down; recoverable failures are
operating-point problems.

\subsection{How do you distinguish laser aging from calibration drift?}
\label{sec:interview-worked-aging-vs-cal}

Separates device physics from control-loop bookkeeping
(\Cref{sec:laser-calibration,sec:laser-thermal-vs-aging}).

\execanswer{Aging changes the LIV. Drift changes the setpoint on a healthy
LIV. External LIV plus recovery after recalibration separates them. Time
behavior confirms: monotonic climb versus a step after a table or firmware
change. Therefore I would route aging to life/derate/replace and drift to
table version control plus an ATP loaded-corner check.}

\paragraph{3-minute answer (practice).}
Remeasure LIV against ship data at fixed junction temperature. Compare
monitor-PD to an external power meter. Sweep modulator bias for a moved
optimum. Aging: life model, derating, burn-in, or replacement. Drift: table
version control, temperature-segment verification, monitor integrity, ATP
corner under load. Do not mix owners.

\paragraph{10-minute reference (read only).}
Monitor-PD corruption is the silent drift mode: APC holds the wrong launch
while telemetry looks fine. A unit that returns to ship performance after a
table reload was never aged.

\subsection{How would you qualify a second laser or photonic-integrated-circuit
supplier?}
\label{sec:interview-worked-second-source}

This question tests supplier judgment, not vendor names. The frame: the first
supplier's failure distribution does not transfer. Qualify against the
requirements slice, not against the incumbent's datasheet
(\Cref{tab:laser-prd,sec:supplier-exec}).

\paragraph{Step 1: freeze the requirements, not the part number.}
Write what the new part must close: link budget, RIN\termnote{RIN} at the
stated ORL\termnote{ORL}, bias window, wavelength class, thermal class, and
lifetime FIT\termnote{FIT} target. Those are the acceptance criteria. The
incumbent's datasheet is evidence that one process can meet them, not a
template the second source must copy. A second source that matches the
datasheet but fails the link-budget corners is not qualified.

\paragraph{Step 2: characterize distributions, not samples.}
Measure threshold, slope, wavelength, SMSR\termnote{SMSR}, and RIN across
wafers, lots, and temperature. Compare spreads against the incumbent, not
only means. A hero sample from a new supplier proves nothing about the
process. Ask for wafer maps and lot genealogy so edge-of-wafer outliers are
visible before they enter your module line. Record the same reference planes
and fixtures you use on the incumbent, or the comparison is fiction.

\paragraph{Step 3: run the full qualification on this process.}
HTOL\termnote{HTOL} with the supplier's own activation-energy justification
for the named mechanism, temperature cycling, damp heat, and burn-in. Wear-
out mechanisms and infant-mortality rates are process-specific: epi, facet
coat, attach, and hermeticity all differ. Borrowing the incumbent's $E_a$ is
the most common quiet mistake. State the projection's fine print: sample
size, confidence bounds, and evidence that the stress accelerates the field
mechanism without inventing a new one.

\paragraph{Step 4: correlate ATPs and plan the ramp.}
Run the same units on the supplier's line and yours so an ATP\termnote{ATP}
limit means the same thing at both sites. Size guard bands from the combined
repeatability. Ramp with lot traceability, a FAIR\termnote{FAIR}, and an
agreed reaction plan for excursions. Keep fleet RMA\termnote{RMA} codes
split by supplier so field data can falsify the qualification. A second
source that cannot be traced in the fleet is not a second source; it is a
blind risk.

\paragraph{How to say it aloud.}
``Requirements first, distributions second, process-specific qual third, ATP
correlation and split field codes fourth.'' Offer to go deep on HTOL validity
or on what you would put in the reaction plan.

\subsection{A single lane is weak in a multi-lane module. How do you isolate
optical, electrical, thermal, and assembly causes?}
\label{sec:interview-worked-lane}

Multi-lane modules give you a free control group: the sibling lanes. The
frame is pattern recognition across lanes before any single-lane deep dive
(\Cref{sec:fm-lane-imbalance}).

\paragraph{Step 1: confirm one outlier, not a gradient.}
Measure every lane for OMA\termnote{OMA}, ER\termnote{ER},
TDECQ\termnote{TDECQ}, and wavelength at the same temperature and pattern. A
single-lane cliff is a local defect. A smooth edge-to-center gradient is
thermal or FAU\termnote{FAU} alignment. A checkerboard pattern often points
at electrical routing or a shared supply. Write the pattern down before
changing anything. The siblings tell you what ``normal'' is for this unit.

\paragraph{Step 2: bisect optical versus electrical.}
Swap the electrical lane assignment, or drive the suspect optical path with
a known-good driver channel. If the weakness follows the optical path, the
candidates are the laser or EML\termnote{EML} element, the modulator, the
attach, or the fiber. If it follows the electrical channel, the candidates
are the driver, the TIA\termnote{TIA}, the package routing, or the host
SerDes\termnote{SerDes} lane. This one swap often cuts the tree in half.

\paragraph{Step 3: separate assembly from device and thermal.}
Per-lane coupled power with a normal LIV\termnote{LIV} on the weak lane
points at fiber-array alignment or a contaminated ferrule: the device is
fine, the light is not leaving. A thermal map or per-lane temperature
monitors separate an edge-lane thermal gradient from a device defect. If the
imbalance grows over HTOL\termnote{HTOL} or burn-in time, one array element
is aging faster, which is a lot or screening question, not a design question.

\paragraph{Step 4: pick the corrective action from the pattern.}
Across lanes, units, and lots, the pattern chooses the fix: a rework
instruction for alignment, a coupling-spec change for the FAU, a thermal-
design change for edge lanes, a driver or TIA lot screen, or a supplier
corrective action on one array element. Do not rewrite the whole module
spec for a single-lane assembly escape.

\paragraph{How to say it aloud.}
``Siblings first, then optical-versus-electrical swap, then assembly versus
device versus thermal, then the pattern picks the control.'' Offer the
FAU-alignment story as the most common production escape.

\subsection{Received power is unchanged but required receiver power increased.
What hypotheses remain?}
\label{sec:interview-worked-sensitivity}

Apply the power-versus-quality fork (\Cref{sec:interview-power-fork}): power
ledger intact, so eye quality or the receiver (\Cref{sec:fm-ber-increase}).

\execanswer{Power held rules out average launch and connector loss as the
primary cause. Remaining hypotheses are Tx eye quality, channel MPI or
wavelength walk, or Rx sensitivity. Therefore I would golden-swap Tx then Rx,
read the BER waterfall shift versus floor, and close on the mechanism and
control.}

\paragraph{3-minute answer (practice).}
List location-tagged hypotheses (Tx ER/RIN/jitter, channel MPI\termnote{MPI}
or filter walk, Rx TIA/PD/equalizer). Golden-swap to own the side. Waterfall
with a VOA\termnote{VOA}: parallel shift means sensitivity moved; a floor
means multiplicative noise. Stressed-eye\termnote{SECQ} if Rx is indicted.
Control follows the mechanism: table, hygiene, supply filter, or lot screen.

\subsection{Why can a link show a BER floor that more launch power does not
fix?}
\label{sec:interview-worked-ber-floor}

This question has a short physics answer and a longer debug answer. Give both
(\Cref{sec:fm-ber-floor,sec:rin}).

\paragraph{Step 1: state the physics in one sentence.}
A BER floor means the dominant noise scales with the signal. Raising launch
power raises that noise in proportion, so the signal-to-noise ratio
saturates. For relative intensity noise the ceiling is
$Q_{\max}=1/\sqrt{\mathrm{RIN}\cdot\mathrm{BW}}$. More power cannot buy you
past that ceiling. Additive noise (thermal, dark current) does not make a
floor; multiplicative noise does.

\paragraph{Step 2: name the usual multiplicative sources.}
Laser RIN, often feedback-driven through poor ORL\termnote{ORL}. Multipath
interference from a pair of reflective interfaces. Bias-rail noise converting
to intensity noise through the laser when PSRR is weak. Crosstalk that
scales with the aggressor. Each of these raises the error rate in a way that
looks like ``the link just will not clean up,'' and each is fixed by
removing the noise source, not by raising OMA.

\paragraph{Step 3: diagnose by waterfall shape and bisect.}
Sweep received power and plot the BER waterfall. Confirm it floors rather
than shifts. Then bisect: measure RIN with a quiet lab source versus the
product bias board to separate intrinsic laser RIN from board-induced
noise. Sweep ORL with a controlled reflector and watch the floor rise and
fall. Read the FEC\termnote{FEC} error histogram: MPI and burst noise cluster
errors in time, while Gaussian noise spreads them. That histogram often
picks the mechanism before a scope does.

\paragraph{Step 4: fix the mechanism, then prevent recurrence.}
Isolator or connector hygiene for feedback. Supply filtering and PSRR
improvement for electrical noise. Aggressor silencing or better isolation for
crosstalk. Replace the laser only if intrinsic RIN is truly out of spec after
the board and ORL have been cleared. Add the ORL or supply stress to the
qualification and to the ATP if that path was never tested.

\paragraph{How to say it aloud.}
``Floor means multiplicative noise; more power cannot help. I confirm with a
waterfall, then bisect RIN, ORL, and the FEC histogram.'' Offer the $Q_{\max}$
line if they want the equation.

\subsection{What makes an HTOL projection credible?}
\label{sec:interview-worked-htol}

A FIT\termnote{FIT} number without a mechanism is a spreadsheet. The frame
is four conditions, said in order (\Cref{sec:laser-aging}).
HTOL\termnote{HTOL} is only as credible as those four.

\paragraph{Condition 1: a named failure mechanism.}
The projection applies to a specific drift parameter: threshold current,
slope efficiency, wavelength, or a defined hard failure such as
COD\termnote{COD}. ``The laser'' is not a mechanism. Name the parameter, the
end-of-life criterion, and the distribution you are projecting.

\paragraph{Condition 2: a justified activation energy.}
$E_a$ must be measured or justified for that mechanism on that process, with
sample size and confidence bounds. Copying an $E_a$ from another product,
another vendor, or a textbook table is the most common way a projection lies.
State the Arrhenius\termnote{Arrhenius} form and the temperatures used, and
be ready to say what happens if $E_a$ is wrong by 0.1~eV.

\paragraph{Condition 3: the stress matches the field mechanism.}
The stress must accelerate the mechanism the fleet will see, without adding a
new one. A stress temperature that triggers solder creep the field never sees
produces a pessimistic fiction. A bench HTOL that omits thermal cycling,
connector wear, or bias-rail transients produces an optimistic fiction. Say
which field mechanisms the HTOL does and does not cover.

\paragraph{Condition 4: a field feedback loop.}
Compare the RMA\termnote{RMA} Pareto against the projection. Divergence is
evidence the model is wrong, not that the fleet is unlucky. Update $E_a$, the
mechanism list, or the use condition when the field disagrees. A projection
that never meets field data is theater.

\paragraph{How to say it aloud.}
``Named mechanism, justified $E_a$, stress matches field, field feedback.
Without those four, FIT is a spreadsheet.'' Offer one example of a stress that
invents a false mechanism.

\subsection{Which data must an automated test save so a failure can be
replayed?}
\label{sec:interview-worked-test-data}

This question is about whether you have ever had to defend a ship decision
months later. The frame: enough data that a unit failing today can be re-
analyzed without the station or the person.

\paragraph{What identity and setup must be saved.}
Sample identity down to serial, lot, and date code. Instrument identity and
calibration state, including the path-loss table version. Fixture, cable, and
reference-plane description, because a power number without a plane is not a
measurement. Software, firmware, and calibration-table versions on the device
under test. Test conditions: temperature, supplies, pattern, dwell, and
timestamps that let results correlate with environmental logs.

\paragraph{What measurement content must be saved.}
Raw measured data, not only pass/fail. Eyes, LIV\termnote{LIV} points,
spectra, and BER\termnote{BER} sweeps that can be re-plotted.
Station-to-station golden-unit correlation results, so measurement drift can
be separated from product drift. Without the raw data, a later engineer can
only argue about the limit, not about the unit.

\paragraph{The replay test of sufficiency.}
Given a field return, can you reproduce the exact shipping measurement and
decide whether the unit changed or the measurement did? If the answer is no,
the archive is incomplete. That question is the acceptance criterion for the
data system, not a nice-to-have.

\paragraph{How to say it aloud.}
``Identity, instrument state, reference plane, raw data, and versions. The
test is whether I can replay ship data on a return.'' Offer one war story
about a missing path-loss table if you have one.

\subsection{When would you choose an EML, silicon MZM, or ring modulator?}
\label{sec:interview-worked-modulator}

Choose by constraint, not preference. State the requirement slice first,
then name what each path buys and costs
(\Cref{tab:source-mod-matrix,tab:tx-modulator}).

\paragraph{When an EML wins.}
An EML\termnote{EML} wins for single-wavelength DR\termnote{DR} or
FR\termnote{FR} at 100--200G/lane when cost, maturity, and one-chip
integration dominate. You get low chirp relative to a DML\termnote{DML}, a
mature supply chain, and a bias-versus-temperature calibration that the
industry already knows how to qualify. You pay for an InP\termnote{InP}
process, EAM\termnote{EAM} aging as a second wear mechanism, and a
calibration table that must be verified at temperature corners.

\paragraph{When a silicon MZM wins.}
A silicon MZM\termnote{MZM} wins when the modulator must sit on the
PIC\termnote{PIC} beside other functions and the link wants a broad, flat
passband with no wavelength lock. You get CMOS\termnote{CMOS}-compatible
integration and a wide optical bandwidth. You pay for die area, drive swing,
RF co-design of the traveling-wave electrode, and a bias-control loop that
holds quadrature. Validation must cover driver matching and bias-rail
behavior across temperature.

\paragraph{When a ring wins.}
A ring\termnote{MRM} wins when many wavelengths must fit on one die, in
dense WDM\termnote{WDM} and co-packaged engines. It is small and WDM-native.
You pay for a wavelength-lock loop, heater power, thermal-crosstalk
validation\termnote{thermalCrosstalk}, and a resonance-alignment failure mode
the others do not have. The validation burden is the decision: if you cannot
afford lock-range and crosstalk testing, you cannot afford the ring.

\paragraph{How to say it aloud.}
``Constraint first. EML for mature single-$\lambda$ DR/FR. Silicon MZM for
broadband PIC integration. Ring for dense WDM with the lock and heater cost
accepted.'' Offer the validation-burden line as the closer.

\subsection{How do you decide whether a field issue is performance,
reliability, or manufacturability?}
\label{sec:interview-worked-triage-class}

Wrong bucket, wrong owner, wasted weeks. Classify on the ticket before
failure analysis starts (\Cref{sec:fleet-triage}).

\paragraph{Bucket 1: performance.}
Ask: did the unit ever meet spec under the field condition? If the design or
operating point never closed the budget at that corner, it is performance.
The fix is retune, derate, a thermal redesign, or a honest spec change.
Shipping more of the same design will not help.

\paragraph{Bucket 2: reliability.}
Ask: did it meet spec at ship and then degrade with time or stress? That is
reliability. The fix is a life model, a screen, derating, burn-in, or field
replacement. Telemetry trends (bias rising at constant power, actuator
heading toward rail) usually show this before the hard failure.

\paragraph{Bucket 3: manufacturability.}
Ask: does the failure cluster by lot, date code, site, or process step
rather than by time? That is manufacturability. The fix is containment,
supplier corrective action (8D\termnote{eightD}), and an ATP\termnote{ATP}
or SPC\termnote{SPC} change. A lot-correlated escape is not a physics
mystery; it is a process-control gap.

\paragraph{How telemetry answers all three before a pull.}
Install age, lot correlation, and trend shape (sudden, gradual, or corner-
dependent) usually pick the bucket before a unit is removed. Classify first,
then pull hardware with a plan that matches the bucket. Performance issues
need a design owner. Reliability issues need a life-model owner.
Manufacturability issues need a supplier and process owner.

\paragraph{How to say it aloud.}
``Ever meet spec? Then degrade with time? Or cluster by lot? Those three
questions pick the owner.'' Offer a one-line example of each bucket from
your experience if you have one.

\subsection{What would you put in fleet telemetry, and why?}
\label{sec:interview-worked-telemetry}

Telemetry exists to catch margin erosion early and to triage without pulling
hardware. Log what discriminates hypotheses, not every register on the chip
(\Cref{sec:cmis,sec:fleet-triage}).

\paragraph{Per-lane observables.}
Transmit and receive optical power, laser bias current, and pre-FEC
BER\termnote{BER} with FEC\termnote{FEC} error histograms. Bias rising at
constant power is aging. Power dropping at constant bias is the optical path.
Clustered errors mean bursts (MPI\termnote{MPI}, connector, ESD event)
rather than Gaussian noise. Per-lane data is what makes the sibling-lane
method possible in the field.

\paragraph{Module observables.}
Temperature, supply rails, TEC\termnote{TEC} or heater drive, and lock-loop
error. Actuators near their rails are consumed margin even while performance
still passes. A railed TEC today is a BER fail next summer. Alarm on actuator
headroom, not only on hard optical thresholds. Read state over
CMIS\termnote{CMIS}.

\paragraph{Events and context.}
LOS\termnote{LOS} and LOL\termnote{LOL} history with timestamps, resets, and
firmware versions, preserved through reboots so intermittents leave evidence.
Lot and date code, install age, and rack position, so one query separates
unit, lot, and environment. Without context, every incident looks unique.

\paragraph{Alarm philosophy.}
Alarm on trends and disagreements, such as monitor versus expected power or
bias versus a peer lane, not only on hard thresholds. Hard thresholds catch
dead units. Trends catch dying ones. The point of fleet telemetry is the
dying ones.

\paragraph{How to say it aloud.}
``Per-lane power, bias, and pre-FEC BER; module temperature and actuator
drive; events with context; alarms on trends.'' Justify each item with the
hypothesis it separates, then stop.

\llmpractice{The worked answers cover laser requirements, the validation
ladder, hot BER with stable power, aging versus calibration, second-source
qual, weak lane isolation, sensitivity shift, BER floor, HTOL credibility,
replayable test data, modulator choice, field triage buckets, and fleet
telemetry. Speak the frame, walk the steps, close on the decision and control.}
{Run a 45-minute mock interview. Pick six of the rehearsal questions at random.
After each answer, ask one follow-up that forces a measurement choice or a
supplier/fleet decision. Score me as Staff-level only if I name the decision
unlocked, not only the instrument used.}

\section{Must-know abbreviations}

A fuller glossary follows this appendix. By the end of the week, know these
without notes. Entries are alphabetical by the leading abbreviation.

\begin{description}
\item[8D / CAPA] Eight-discipline problem solving and corrective and preventive
      action. Structured containment, root cause, correction, and recurrence
      control for supplier or production failures.
\item[APC] Automatic power control. Feedback from a monitor photodiode that
      holds average launch power; a drifting monitor corrupts it silently.
\item[Arrhenius / $E_a$] Life-acceleration model. Temperature stress multiplies
      wear-out by $\exp[(E_a/k)(1/T_\mathrm{use}-1/T_\mathrm{stress})]$. $E_a$
      must be justified for the named mechanism on that process.
\item[ATP] Acceptance test plan. Production test limits, methods, reference
      planes, and reaction rules.
\item[Bathtub curve] Infant mortality (falling rate), useful life (roughly
      constant FIT), then wear-out (rising rate). Burn-in targets the left;
      Arrhenius life targets the right.
\item[BER] Bit error ratio. Pre-FEC BER is the validation metric; post-FEC is
      the service metric. Always say which.
\item[BERT] Bit-error-ratio tester. Generates PRBS and counts errors for
      waterfalls, floors, and dwell tests.
\item[Burn-in] Production or sample screen that removes infant-mortality parts
      before ship. Distinct from HTOL life projection.
\item[CDR] Clock-data recovery. Extracts bit clock from the data stream; LOL
      means unlock even if light is present.
\item[CMIS] Common Management Interface Specification. Module state, controls,
      alarms, and telemetry.
\item[COD] Catastrophic optical damage. Sudden, irreversible laser-facet failure.
\item[COM] Channel operating margin. Statistical electrical-link margin after
      loss, noise, crosstalk, and equalization.
\item[DCA] Digital communication analyzer. Sampling oscilloscope for eyes,
      OMA, ER, RLM, and TDECQ.
\item[DFB] Distributed-feedback laser. Single-mode 1310~nm class source; also
      the gain section in an EML.
\item[DML] Directly modulated laser. Simple and efficient, but chirp limits reach.
\item[DPA] Destructive physical analysis. Cross-section or EDX on failed units
      to confirm facet, solder, FAU, or die failure modes.
\item[DPPM] Defective parts per million. Incoming or outgoing quality rate.
\item[DR / FR] Datacenter reach ($\sim$500~m) and far reach (2~km). IEEE
      single-mode classes at 1310~nm.
\item[DVT / PVT] Design validation test and production validation test.
\item[EOL] End of life. The defined wear-out criterion (threshold rise, slope
      drop, or hard fail) used in HTOL projection and derating.
\item[ESD] Electrostatic discharge. Handling or assembly damage to drivers or
      TIAs; sudden hard fail, not Arrhenius wear-out. Qual uses HBM/CDM models.
\item[EAM] Electro-absorption modulator. Voltage-controlled absorption; pairs
      with a DFB in an EML.
\item[ELSFP] External Laser Small Form-Factor Pluggable. Replaceable continuous-
      wave source for co-packaged optics.
\item[EML] Electro-absorption modulated laser. DFB plus EAM on one InP chip;
      dominant 100--200G/lane pluggable transmitter.
\item[ER] Extinction ratio. $P_1/P_0$ in dB. Higher ER widens OMA at fixed
      average power; trades against chirp and swing.
\item[FAIR] First-article inspection report. Re-qualify after tooling, site,
      epi, or firmware change before open volume.
\item[FAU] Fiber array unit. Precision multi-fiber attachment to a photonic
      integrated circuit.
\item[FEC / KP4] Forward error correction; KP4 is Reed--Solomon RS(544,514).
      Pre-FEC BER and error histograms show margin before FEC failure.
\item[FIT] Failures in time. Failures per $10^9$ device-hours.
\item[Golden unit / host] Known-good reference used to bisect station, fixture,
      host, module, and fiber. Essential in debug and ATP correlation.
\item[GR-468] Telcordia optoelectronic qualification framework (HTOL,
      environmental, mechanical).
\item[HAST] Highly accelerated stress test. Humidity plus temperature and bias;
      exercises corrosion and delamination.
\item[HTOL] High-temperature operating life. Accelerated stress used with an
      Arrhenius model to project field wear-out; credible only with a named
      mechanism.
\item[HTSL] High-temperature storage life. Unbiased bake; separates storage
      mechanisms from biased HTOL wear-out.
\item[JESD47] JEDEC IC qualification stress suite. Silicon-side counterpart to
      GR-468 for drivers and TIAs.
\item[LIV] Light--current--voltage curve. Threshold, slope efficiency, kink-free
      range, and thermal rollover versus bias.
\item[LOS / LOL] Loss of signal and loss of lock. LOS points first toward power;
      LOL can occur with adequate power but poor timing or signal quality.
\item[MPI] Multipath interference. Coherent beating from reflective paths;
      often floors BER and clusters FEC errors.
\item[MPO] Multi-fiber push-on connector. Parallel-optics plant (8--32 fibers).
\item[MRM] Microring modulator. Compact, WDM-native silicon modulator; needs
      wavelength lock and heater budget.
\item[MSA] Multi-source agreement. An industry specification for interoperable
      products.
\item[MTBF] Mean time between failures. For constant failure rate,
      $\mathrm{MTBF}=10^9/\mathrm{FIT}$ hours. Fleet math usually uses FIT
      times population.
\item[MZM] Mach--Zehnder modulator. Broadband interferometric modulator (Si or
      TFLN); needs quadrature bias control.
\item[NFF / RMA] No fault found and return merchandise authorization. High NFF
      rates often indicate weak triage or intermittent faults.
\item[OMA] Optical modulation amplitude. Outer level swing $P_1-P_0$.
\item[ORL] Optical return loss. Reflected power toward the laser; low ORL raises
      RIN and can seed burst errors.
\item[OSA] Optical spectrum analyzer. Wavelength, SMSR, and side-mode structure.
\item[OSFP / QSFP-DD] High-density pluggable module form factors with different
      mechanical and thermal limits.
\item[PIC / SOI] Photonic integrated circuit on silicon-on-insulator. The common
      silicon-photonics chip and substrate.
\item[PRBS] Pseudo-random binary sequence. Repeatable pattern for eye and BER
      measurements.
\item[PSRR] Power-supply rejection ratio. Weak PSRR can turn electrical rail
      noise into optical intensity noise.
\item[RIN] Relative intensity noise. Laser amplitude noise (dB/Hz); often
      measured under a stated ORL.
\item[RLM] Relative level mismatch. PAM4 level-spacing quality.
\item[SECQ] Stressed eye closure quaternary. Receiver-side margin measured with a
      calibrated stressed optical signal.
\item[SerDes] Serializer/deserializer. Host high-speed I/O; equalization reserve
      is a timing-margin ledger.
\item[SMSR] Side-mode suppression ratio. Power difference between the lasing mode
      and the strongest side mode.
\item[SPC] Statistical process control. Tracks process distributions and trends.
\item[TDECQ] Transmitter and dispersion eye closure quaternary. Headline PAM4
      transmitter-quality metric after a reference receiver and bounded FFE.
\item[TEC] Thermoelectric cooler. A TEC near its current limit has little thermal
      control margin left.
\item[TIA] Transimpedance amplifier. Converts photodiode current to voltage.
\item[VCSEL] Vertical-cavity surface-emitting laser. 850~nm class for multimode
      short-reach links.
\item[VNA] Vector network analyzer. Electrical or electro-optic $S$-parameters
      versus frequency.
\item[VOA] Variable optical attenuator. Calibrated loss for sensitivity and
      BER-waterfall sweeps.
\item[WDM] Wavelength division multiplexing. Multiple wavelengths on one fiber.
\end{description}

\llmpractice{Know the alphabetical list cold. Optical/debug core: ATP, APC,
BER/BERT, CMIS, DCA, EML, ER, FEC/KP4, FIT, HTOL, LIV, LOS/LOL, MPI, OMA,
ORL, RIN, RLM, SECQ, TDECQ, TEC, TIA, VOA. Reliability/manufacturing core:
Arrhenius/$E_a$, bathtub, burn-in, DPA, ESD, FAIR, golden unit, GR-468, HAST,
HTSL, JESD47, MTBF, NFF/RMA, SPC, 8D/CAPA. Expand the term, then say why it
matters in a debug or ship decision.}
{Drill abbreviations. Give me ten random terms mixing optical debug and
reliability/manufacturing. For each, I must expand it in one sentence and give
one measurement, stress, or failure mode it connects to. Fail me if I only
expand the letters or confuse burn-in with HTOL.}

\section{One-week study plan}
\label{sec:interview-week-plan}

Assume seven days, with the interview near the end of Day~7. Protect sleep.
Each day has one primary drill and one LLM practice box. If a day slips, cut
new reading first, not story rehearsal or the mock interview.

\paragraph{Day 1: principles, decision trees, spine, traps, Staff judgment.}
Read the four principles, \Cref{sec:interview-decision-trees}, the answer
spine, the traps, and \Cref{sec:interview-under-uncertainty}. Memorize: every
answer ends with a decision; measurements characterize margin; debugging is
Bayesian inference in the lab. Memorize the two universal trees, the four
phases, the six-question checklist, the decision table, and the five ledgers.
Speak one paragraph per phase. Use the answer-spine LLM practice box. Write
ownership language: stop ship, notify supplier, request FA, update ATP,
monitor RMA.

\paragraph{Day 2: requirements and validation ladder.}
Rehearse \Cref{sec:interview-worked-laser-reqs} and
\Cref{sec:interview-worked-validation} aloud until the structure is automatic.
Name a reference plane in every measurement sentence. Skim
\Cref{tab:ladder,tab:laser-prd,tab:source-mod-matrix} for numbers you already
believe; do not hunt new optics.

\paragraph{Day 3: two real stories.}
Draft and rehearse one bench or component story and one system, production, or
fleet story. Use only work you did. Hit all eight beats and end on a measured
recurrence control. Use the stories LLM practice box. Record yourself once and
cut invented metrics.

\paragraph{Day 4: instruments, waterfall, and debug forks.}
Review what each instrument removes as uncertainty. Drill waterfall shift
versus floor versus burst pattern
(\Cref{sec:interview-waterfall,sec:fm-ber-increase,sec:fm-ber-floor}). Rehearse
hot BER with stable power, aging versus calibration, and weak-lane isolation.
Use the ten-concepts LLM practice box.

\paragraph{Day 5: reliability, manufacturing, and suppliers.}
Drill HTOL credibility, Arrhenius/$E_a$, bathtub (burn-in versus wear-out),
second-source qualification, FAIR, DPA, ESD versus wear-out, and field triage
buckets (performance / reliability / manufacturability). Rehearse the HTOL,
second-source, and triage worked answers. Skim \Cref{ch:reliability} only where
a story needs a fact.

\paragraph{Day 6: abbreviations, telemetry, and modulator choice.}
Run the abbreviations LLM practice box until expansions are fast. Rehearse
fleet telemetry, modulator choice (EML / Si MZM / ring), replayable test data,
and BER-floor physics. Do one short mixed quiz: three debug questions and three
reliability questions.

\paragraph{Day 7: cheat sheet, Appendix~B frameworks, mock, and stop.}
Read \Cref{sec:interview-cheat-sheet} once aloud. Drill the green
thirty-second boxes in \Cref{app:interview-frameworks} for the topics you
expect. Use the 45-minute mock-interview LLM practice box. No new chapters
beyond that drill. Light review of your two stories only. Stop two to three
hours before the call. Sleep.

\paragraph{If you have less than seven days.}
Compress in this order: Day~1 spine, Day~3 stories, Day~2 requirements and
validation, Day~5 HTOL and triage, Day~7 mock. Cut Day~6 reading; keep only the
abbreviation drill. Always keep the cheat sheet below.

\clearpage
\section{Staff interview cheat sheet}
\label{sec:interview-cheat-sheet}

Internalize this one page. The rest of the appendix is supporting detail.

{\small\setlength{\parskip}{0.3em}%
\noindent\textit{Close.} Every answer ends with the engineering decision.\par
\noindent\textit{Philosophy.} Engineering reduces uncertainty. Measurements
unlock decisions. Measurements characterize margin. Every measurement updates
belief. The job is the best decision with today's evidence.\par
\noindent\textit{Trees.} Debug: scope $\rightarrow$ power fork $\rightarrow$
isolation $\rightarrow$ decision $\rightarrow$ control. Qual: works $\rightarrow$
margins $\rightarrow$ environment $\rightarrow$ interop $\rightarrow$ life
$\rightarrow$ ATP $\rightarrow$ fleet (\Cref{sec:interview-decision-trees}).\par
\noindent\textit{Spine (4 phases).} Understand $\rightarrow$ Investigate
$\rightarrow$ Resolve $\rightarrow$ Prevent.\par
\noindent\textit{Checklist.} Scope? Ledger? Power or quality? Fastest
measurement? Decision? Control?\par
\noindent\textit{Night before.} Open the matching framework in
\Cref{app:interview-frameworks}. Memorize green 30-second boxes only.\par
\noindent\textit{Five ledgers.} Power $\cdot$ Noise $\cdot$ Timing $\cdot$
Spectral $\cdot$ Control.\par
\noindent\textit{Margin budget.} Stress consumes margin; qual measures what
remains (\Cref{sec:interview-margin-budget}).\par
\noindent\textit{Customer vs vendor.} Customer measures external behavior;
vendor owns internals (\Cref{sec:interview-customer-vendor}).\par
\noindent\textit{Decisions.} Ship / don't ship, continue validation, escalate
supplier, derate, second source, contain lot, modify ATP, open RMA, FA/DPA,
firmware, retune calibration, monitor only (\Cref{tab:interview-decisions}).\par
\noindent\textit{Traps / ownership.} No component-first; no stop after fix.
Contain, notify, FA, ATP, monitor RMA.\par
}

\keyidea{I first want to understand the scope of the problem, then determine
which margin ledger is being spent, choose the measurement that eliminates the
largest number of hypotheses, make the product decision, and finally add the
control that prevents the next escape.}
